\labsection{5}{Neobank unit economics}{sec:lab05-unit-economics}

\begin{labproject}{Find out whether a digital-only bank makes money per customer}
\textbf{Coverage.} Chapter~\ref{ch:digital-banks}.\\
\textbf{Goal.} Build the model that answers the chapter's third strategic priority --- expand \emph{profitably} --- and find the assumption the whole business rests on.

\courselabflow{Set assumptions}{Compute margin}{Project cash}{Vary inputs}



\end{labproject}

\textbf{Complete worked example.} This runnable fixture uses synthetic inputs. First reproduce it, then adapt the data or ratings with documented evidence. All code below belongs to one file.
\begin{lstlisting}[style=utpterminal]
python all_experiments/lab05_neobank_unit_economics.py
\end{lstlisting}

\noindent\textbf{Setup and output folder.}
\lstinputlisting[style=utppython,firstline=1,lastline=10]{all_experiments/lab05_neobank_unit_economics.py}

\noindent\textbf{Define and validate unit economics.}
\lstinputlisting[style=utppython,firstline=11,lastline=24]{all_experiments/lab05_neobank_unit_economics.py}

\noindent\textbf{Project the customer stock and cumulative cash.}
\courselabtrace{Profitable contribution does not mean positive cash}{Contribution}{4.20 - 1.80 = 2.40 per customer-month}{Acquisition}{CAC 45 is paid when each new customer arrives}{Payback}{45 / 2.40 = 18.75 months before churn or fixed costs}

\textbf{Use the chapter while coding.} Revisit \hyperref[ch:digital-banks]{the earlier chapter explanation}; match its operation to the variables and state changes below.

\textbf{Follow the supplied implementation.} In project, churn applies to existing customers before new customers are added. New customers earn a full month of contribution under this teaching convention. cash += accumulates margin minus acquisition spending, so the returned cash is a balance rather than monthly profit. Compare the first two rows of fast.csv with this trace before interpreting the 60-month chart. Values are synthetic and omit fixed costs and financing.
\courselabtrace{Trace the fast-acquisition cash projection}{1. Input / state}{Month 1: 50,000 new customers. Contribution 120,000 minus acquisition cost 2,250,000 gives cash -2,130,000.}{2. Operation}{Month 2: 50,000 times 0.97 + 50,000 = 98,500 customers. Monthly contribution = 236,400.}{3. Check}{Month 2 cash = -2,130,000 + 236,400 - 2,250,000 = -4,143,600. Cash is cumulative.}


\lstinputlisting[style=utppython,firstline=25,lastline=46]{all_experiments/lab05_neobank_unit_economics.py}

\noindent\textbf{Change one assumption at a time.}
\lstinputlisting[style=utppython,firstline=47,lastline=51]{all_experiments/lab05_neobank_unit_economics.py}
\begin{labvisual}{Synthetic projection: contribution versus acquisition cash}
\begin{labgraphic}\begin{tikzpicture}
\begin{axis}[width=12cm,height=5.3cm,xlabel={Month},ylabel={Cumulative cash (millions)},
tick label style={font=\scriptsize},label style={font=\small},grid=major,
legend style={font=\scriptsize,at={(.03,.97)},anchor=north west},xmin=1,xmax=60]
\addplot[thick,color=UTPBlue] coordinates {(1,-2.13)(2,-4.1436)(3,-6.044292)(4,-7.83546324)(5,-9.520399342800001)(6,-11.102287362516002)(7,-12.584218741640521)(8,-13.969192179391307)(9,-15.260116414009568)(10,-16.459812921589283)(11,-17.571018533941604)(12,-18.596387977923357)(13,-19.538496338585656)(14,-20.399841448428088)(15,-21.182846204975245)(16,-21.889860818825987)(17,-22.523164994261204)(18,-23.08497004443337)(19,-23.57742094310037)(20,-24.002598314807358)(21,-24.36252036536314)(22,-24.659144754402245)(23,-24.89437041177018)(24,-25.070039299417076)(25,-25.187938120434563)(26,-25.249799976821528)(27,-25.257305977516882)(28,-25.212086798191375)(29,-25.115724194245637)(30,-24.969752468418267)(31,-24.77565989436572)(32,-24.53489009753475)(33,-24.24884339460871)(34,-23.91887809277045)(35,-23.54631174998734)(36,-23.13242239748772)(37,-22.67844972556309)(38,-22.185596233796197)(39,-21.655028346782313)(40,-21.087877496378848)(41,-20.485241171487484)(42,-19.84818393634286)(43,-19.17773841825258)(44,-18.474906265705005)(45,-17.740659077733856)(46,-16.975939305401845)(47,-16.18166112623979)(48,-15.358711292452599)(49,-14.507949953679024)(50,-13.630211455068656)(51,-12.726305111416599)(52,-11.797015958074104)(53,-10.843105479331884)(54,-9.86531231495193)(55,-8.864352945503375)(56,-7.840922357138276)(57,-6.795694686424131)(58,-5.729323845831409)(59,-4.64244413045647)(60,-3.535670806542779)};
\addlegendentry{Continue acquisition}
\addplot[thick,color=UTPOrange] coordinates {(1,-2.13)(2,-4.1436)(3,-6.044292)(4,-7.83546324)(5,-9.520399342800001)(6,-11.102287362516002)(7,-12.584218741640521)(8,-13.969192179391307)(9,-15.260116414009568)(10,-16.459812921589283)(11,-17.571018533941604)(12,-18.596387977923357)(13,-19.538496338585656)(14,-20.399841448428088)(15,-21.182846204975245)(16,-21.889860818825987)(17,-22.523164994261204)(18,-23.08497004443337)(19,-23.57742094310037)(20,-24.002598314807358)(21,-24.36252036536314)(22,-24.659144754402245)(23,-24.89437041177018)(24,-25.070039299417076)(25,-23.057938120434564)(26,-21.106199976821525)(27,-19.213013977516884)(28,-17.376623558191376)(29,-15.595324851445637)(30,-13.86746510590227)(31,-12.1914411527252)(32,-10.565697918143446)(33,-8.988726980599143)(34,-7.4590651711811695)(35,-5.975293216045736)(36,-4.5360344195643645)(37,-3.1399533869774348)(38,-1.7857547853681124)(39,-0.47218214180706974)(40,0.8019833224471414)(41,2.037923822773726)(42,3.2367861080905134)(43,4.3996825248477975)(44,5.527692049102362)(45,6.62186128762929)(46,7.683205449000411)(47,8.712709285530398)(48,9.711328006964484)(49,10.679988166755548)(50,11.619588521752881)(51,12.531000866100293)(52,13.415070840117282)(53,14.272618714913763)(54,15.10444015346635)(55,15.911306948862357)(56,16.693967740396484)(57,17.45314870818459)(58,18.18955424693905)(59,18.903867619530875)(60,19.596751590944947)};
\addlegendentry{Stop after month 24}
\end{axis}\end{tikzpicture}\end{labgraphic}\end{labvisual}

\textbf{Check the run.} The following output was captured from the supplied script; it describes the synthetic fixture only.
\lstinputlisting[style=utpterminal]{labs/outputs/lab05.txt}
% Reused from chapters/05-digital-only-banks.tex
\begin{labvisual}{Recall the chapter model}
\begin{labgraphic}
\begin{tikzpicture}[
  b/.style={draw=UTPBlue,fill=UTPPaleBlue,rounded corners=2mm,text width=2.65cm,minimum height=1.25cm,align=center,font=\scriptsize\bfseries\color{UTPNavy}},
  g/.style={draw=UTPGreen,fill=UTPPaleTeal,rounded corners=2mm,text width=2.65cm,minimum height=1.25cm,align=center,font=\scriptsize\bfseries\color{UTPNavy}},
  arr/.style={-{Stealth[length=2mm]},thick,draw=UTPTeal}
]
\node[b] (a) at (-4.5,0) {Modern modular\\technology};
\node[b] (b) at (-1.5,0) {Lower servicing\\cost};
\node[b] (c) at (1.5,0) {Fast product and\\CX iteration};
\node[g] (d) at (4.5,0) {Scale + customer\\engagement};
\draw[arr](a)--(b); \draw[arr](b)--(c); \draw[arr](c)--(d);
\end{tikzpicture}
\end{labgraphic}\end{labvisual}


\begin{verification}
\textbf{Explain and extend your result.} Change churn while holding other assumptions fixed. Compare LTV/CAC and the lowest monthly cash balance; explain why they answer different business questions.

\begin{itemize}\item The script runs from a fresh process and creates the named CSV or report files.\item Every rate uses a stated population and denominator; every model assumption is identified.\item Synthetic outputs are labelled and are not presented as observations of a real institution.\item Change one input, rerun, and explain the direction of the result.\end{itemize}\end{verification}
\begin{labdeliverable}Submit the complete script, generated evidence, and a one-page interpretation linking the results to the chapter model. State the assumptions and one limitation.\end{labdeliverable}
